\documentclass[10pt, conference, letterpaper]{IEEEtran}

\usepackage{cite}
\usepackage{amsmath,amssymb,amsfonts}
\usepackage{graphicx}
\usepackage{booktabs}
\usepackage{url}
\usepackage{float} 

\begin{document}

\title{Prompt Sensitivity and Its Effect on the Stability of LLM-Generated Research Conclusions: An Empirical Implementation\\
}

\author{
\IEEEauthorblockN{Dhyan Shah}
\IEEEauthorblockA{\textit{M.Tech (IT), Semester I} \\
\textit{IIIT Allahabad}\\
Prayagraj, Uttar Pradesh, India}
}

\maketitle

\begin{abstract}
Large Language Models (LLMs) are increasingly deployed in multi-agent workflows to synthesize research and extract empirical conclusions. However, single-turn prompts are notoriously sensitive to minor perturbations in formatting, semantics, and candidate positioning. This paper translates the theoretical vulnerabilities identified in prompt sensitivity literature into an executable Monte Carlo simulation of a 4-stage agentic workflow. Utilizing the PubMedQA dataset and two open-weight models (Qwen2.5-0.5B-Instruct and Qwen2.5-1.5B-Instruct), we systematically evaluate 8 meaning-preserving prompt variants. Our empirical results demonstrate that while stability improves with parameter scale, both models remain highly fragile. Specifically, surface formatting changes induce a 100\% conclusion flip rate in the 0.5B model, while the 1.5B model exhibits severe vulnerability to positional option reversal (46.7\% flip rate) and semantic paraphrasing (50.0\% flip rate). Furthermore, comparing prompt-variant agreement against a repeat-control baseline confirms that the observed instability is structurally driven by the prompt string rather than sampling variance. We conclude that prompt variations cause severe semantic claim drift, and that scale attenuates, but does not eliminate, the risk of cascading failures in multi-stage pipelines.
\end{abstract}

\begin{IEEEkeywords}
prompt sensitivity, large language models, conclusion stability, cascading failures, multi-agent workflows
\end{IEEEkeywords}

\section{Introduction}
Researchers are rapidly integrating Large Language Models (LLMs) into complex, multi-stage pipelines to automate literature synthesis, evidence classification, and comparative method analysis \cite{wang2024survey, xi2023rise}. While these autonomous tools offer significant efficiency gains, their fundamental reliability remains a critical concern \cite{mizrahi2024state}. A fluent, seemingly authoritative response generated by an LLM is heavily dependent on the specific instruction sequence provided in the prompt \cite{reynolds2021prompt, zhou2023large}. 

Prior literature has firmly established that single-turn LLM outputs are highly sensitive to superficial input variations \cite{sclar2024quantifying, salinas2024}. A minor change in whitespace, template formatting, or the order of few-shot examples can drastically alter a model's prediction \cite{lu2022fantastically}. However, in the context of autonomous research assistants, these isolated errors do not occur in a vacuum. A hallucination or formatting breakdown in an early stage (such as document retrieval) serves as the corrupted input state for subsequent downstream stages (such as evidence extraction and verification) \cite{mialon2023augmented, wu2023agentbench}. 

Despite this known fragility, most evaluations of prompt sensitivity still rely on benchmark accuracy or simple classification labels \cite{brown2020language, touvron2023llama}. In automated research workflows, the final output is typically an open-ended scientific claim. Two generated responses may feature high lexical similarity but differ in critical qualifiers, effectively reversing the extracted conclusion \cite{elazar2021measuring}. Consequently, measuring textual variance or benchmark accuracy is insufficient for evaluating the reliability of LLM-assisted scientific research \cite{turpin2023language}.

This paper bridges the critical gap between single-turn prompt sensitivity and multi-step pipeline stability. We investigate the following central research question: \textit{How do compounding error probabilities driven by prompt sensitivity affect the terminal failure rate and stability of research conclusions in LLM-assisted workflows, compared to an independent error baseline?}

Through a controlled comparative study, we successfully separate standard sampling variance (thermal decoding noise) from prompt-driven structural instability \cite{wang2023self}. By defining a Conclusion Agreement Rate (CAR) and applying normalized entropy metrics, we isolate changes in the actual scientific claim from mere variations in wording \cite{errica2025wrong}. Our findings highlight the absolute necessity of implementing rigorous prompt variant testing and ensemble verification architectures before accepting LLM-generated research claims \cite{zhao2021calibrate}.

\section{Theoretical Background and Related Work}
Our methodology builds directly upon extensive prior literature in prompt engineering robustness, cognitive biases, and agentic reliability.

\subsection{Surface Formatting Sensitivity}
Recent studies have demonstrated that modern LLMs exhibit severe vulnerability to seemingly cosmetic formatting changes \cite{sclar2024quantifying, salinas2024}. Sclar et al. \cite{sclar2024quantifying} conducted extensive evaluations on few-shot classification settings, revealing that altering separators, changing whitespace, or shifting the requested output structure caused massive performance drops in models like LLaMA-2. This suggests that models heavily overfit to the formatting distributions seen during instruction-tuning \cite{ouyang2022training}, and deviations from these templates disrupt their internal reasoning pathways \cite{huang2023towards}.

\subsection{Semantic Sensitivity and Consistency}
While formatting sensitivity addresses the structure of the prompt, semantic sensitivity addresses the phrasing of the instruction itself \cite{webson2022do}. Errica et al. \cite{errica2025wrong} introduced rigorous mathematical frameworks to separate a model's true task robustness from its baseline performance. By utilizing meaning-preserving paraphrases, they defined Sensitivity (via normalized entropy) and Consistency (via Total Variation Distance) to quantify how deeply semantic rewording degrades model outputs \cite{elazar2021measuring}. 

\subsection{Positional Bias in LLM-as-a-Judge Architectures}
When models are used to evaluate, extract, or rank alternatives, the order of the candidates introduces significant bias \cite{lu2022fantastically}. Zheng et al. \cite{zheng2023judging} identified position, verbosity, and self-enhancement biases in LLM-as-a-judge architectures. In a research synthesis context, this implies that the order in which retrieved evidence or candidate conclusions are concatenated into the prompt can artificially skew the final extracted claim.

\section{Experimental Methodology}
To empirically test the stability of LLM-generated conclusions under controlled perturbations, we engineered a reproducible Phase 3 Python pipeline executed on a Google Colab environment utilizing a Tesla T4 GPU. 

\subsection{Dataset and Sampling Strategy}
We utilized the \texttt{qiaojin/PubMedQA} dataset (specifically the pqa\_labeled subset) \cite{jin2019pubmedqa}. PubMedQA is an expert-annotated biomedical question-answering dataset where the task is to answer a research question based on a provided abstract with a `yes', `no', or `maybe' conclusion. We deterministically sampled 10 instances from each of the three gold-standard conclusion classes, resulting in 30 total evaluation questions \cite{honovich2022true}. 

\subsection{Models and Decoding Parameters}
We evaluated two open-weight, instruction-tuned models from the Qwen family to analyze the effect of parameter scale on stability: \texttt{Qwen2.5-0.5B-Instruct} and \texttt{Qwen2.5-1.5B-Instruct} \cite{qwen2024}. For the primary perturbation evaluation, we utilized deterministic (greedy) decoding (\texttt{do\_sample=False}) to ensure that any observed variance across outputs is strictly the result of the prompt perturbation, rather than thermal sampling noise \cite{kojima2022large}.

\subsection{Prompt Perturbation Taxonomy}
We defined a baseline prompt ($P_0$) requesting a neutral yes/no/maybe conclusion and a concise rationale. We then engineered 7 meaning-preserving variants mapped across four specific perturbation classes:
\begin{itemize}
    \item \textbf{Baseline ($P_0$):} Standard neutral instruction.
    \item \textbf{Surface/Formatting ($P_1, P_2$):} Altered delimiters to use capitalization and hash symbols, and requested the output in a strict JSON format \cite{sclar2024quantifying}.
    \item \textbf{Semantic Paraphrase ($P_3, P_4$):} Highly verbose academic rewording and concise clinical synthesis rewording \cite{errica2025wrong}.
    \item \textbf{Positional/Order ($P_5$):} Reversed the acceptable option string from ``yes, no, or maybe'' to ``maybe, no, or yes'', and inverted the order of the Abstract and Question \cite{lu2022fantastically}.
    \item \textbf{Contextual/Framing ($P_6, P_7$):} Applied a polite confirmatory peer-reviewer persona and a strict, skeptical auditor persona \cite{webson2022do}.
\end{itemize}

To rigorously isolate prompt-driven instability from inherent generative variance, we also executed 5 independent repeat runs of the baseline prompt using non-deterministic decoding (\texttt{temperature=0.3}) \cite{wang2023self}. This serves as the ``Repeat Control'' baseline.

\subsection{Metrics and Claim Extraction}
To measure stability at the level of the scientific claim rather than raw text, we implemented a deterministic claim-extraction parser, $C(y)$ \cite{ji2023survey}. The stability of the parsed conclusions was evaluated using three primary metrics:
1) \textit{Conclusion Agreement Rate (CAR):} Calculates the relative frequency of the modal (most common) conclusion across a set of variants \cite{mizrahi2024state}.
2) \textit{Errica Sensitivity ($S_\tau$):} Adopts normalized entropy to measure the label distribution across prompt variants \cite{errica2025wrong}.
3) \textit{Mean Wording Similarity:} Evaluates pairwise lexical similarity to separate surface wording variation from actual claim agreement \cite{kamalloo2023evaluating}.

\section{Results and Deeper Analysis}

The evaluation generated 450 total responses per model. The overarching summary of model stability is presented in Table \ref{tab:summary}, and the specific perturbation flip rates are detailed in Table \ref{tab:breakdown}.

\begin{table*}[htbp]
\centering
\caption{Model Comparison Summary: Accuracies, Agreement Rates, and Sensitivity Metrics}
\label{tab:summary}
\begin{tabular}{@{}lcccccccc@{}}
\toprule
\textbf{Model} & \textbf{P0 Acc.} & \textbf{Modal Acc.} & \textbf{CAR (PV)} & \textbf{CAR (RC)} & \textbf{Word Sim.} & \textbf{Sens. ($S_\tau$)} & \textbf{Consistency} & \textbf{Unparsed} \\ \midrule
0.5B-Instruct & 0.300 & 0.367 & 0.500 & 0.913 & 0.202 & 0.847 & 0.758 & 0.013 \\
1.5B-Instruct & 0.500 & 0.467 & 0.750 & 0.813 & 0.166 & 0.471 & 0.623 & 0.008 \\ \bottomrule
\end{tabular}
\end{table*}

\subsection{Scale Attenuates, but Does Not Eliminate, Instability}
The aggregated data demonstrates a clear trend: stability improves with parameter scale, yet remains unacceptably fragile for autonomous research applications \cite{qwen2024}. The 1.5B model achieved a multi-prompt Conclusion Agreement Rate (CAR) of 0.750 and a normalized entropy sensitivity score ($S_\tau$) of 0.471, substantially outperforming the 0.5B model (CAR 0.500, $S_\tau$ 0.847). 

However, a CAR of 0.750 in the larger model still implies that 25\% of the time, the model contradicts its own modal conclusion when presented with a meaning-preserving prompt variant \cite{errica2025wrong}. In a multi-stage pipeline, an inherent 25\% error rate at the extraction stage practically guarantees cascading failures in downstream verification stages \cite{wu2023agentbench}.

\subsection{Catastrophic Failures on Surface Formatting (0.5B)}
Analysis of the 0.5B-Instruct model reveals severe token overfitting to instructional templates \cite{sclar2024quantifying}. Merely altering the prompt separators resulted in a 100\% conclusion flip rate compared to the baseline prompt. Furthermore, requesting a JSON format resulted in a 93.3\% flip rate. These catastrophic failures suggest that lower-parameter models heavily memorize specific Markdown layouts during instruction-tuning \cite{ouyang2022training}. When a prompt deviates from these layouts, the model's internal reasoning trace is completely disrupted.

\begin{table}[htbp]
\centering
\caption{Perturbation Breakdown (Conclusion Flip Rates vs P0 Baseline)}
\label{tab:breakdown}
\begin{tabular}{@{}llcc@{}}
\toprule
\textbf{Category} & \textbf{Variant ID} & \textbf{0.5B Flip} & \textbf{1.5B Flip} \\ \midrule
Surface & P1\_separators & 1.000 & 0.333 \\
Surface & P2\_json\_format & 0.933 & 0.167 \\
Semantic & P3\_para\_1 & 0.967 & 0.367 \\
Semantic & P4\_para\_2 & 0.033 & 0.500 \\
Positional & P5\_reverse\_opt & 0.967 & 0.467 \\
Contextual & P6\_confirm & 0.667 & 0.300 \\
Contextual & P7\_skeptical & 0.033 & 0.300 \\ \bottomrule
\end{tabular}
\end{table}

\subsection{Positional and Semantic Vulnerability (1.5B)}
Interestingly, while scaling to 1.5 billion parameters effectively mitigated the surface formatting fragility (reducing the flip rates significantly), the 1.5B model revealed deep vulnerabilities to structural and semantic perturbations \cite{qwen2024}.

Reversing the string order of the acceptable options from ``yes, no, or maybe'' to ``maybe, no, or yes'' caused a massive 46.7\% flip rate. This empirically confirms the positional bias theories proposed by Zheng et al. \cite{zheng2023judging}: the model anchors heavily on the sequence of tokens presented in the prompt rather than grounding its choice strictly in the provided biomedical abstract \cite{lu2022fantastically}.

Moreover, a concise clinical paraphrase triggered a 50.0\% flip rate \cite{errica2025wrong}. This proves that researchers interacting with the 1.5B model have a literal coin-flip chance of receiving contradictory claims simply by rewording their question, invalidating the model's reliability as a standalone extraction agent.

\begin{figure*}[htbp]
    \centering
    \includegraphics[width=0.8\textwidth]{1.1.png}
    \caption{Conclusion Instability (Flip Rates relative to the P0 Baseline) categorized by perturbation class. The 0.5B model fails catastrophically on surface formatting, while the 1.5B model remains highly vulnerable to semantic and positional shifts.}
    \label{fig:flip_rates}
\end{figure*}

\subsection{Disentangling Prompt Variance from Sampling Variance}
A critical requirement of our methodology was proving that the observed instability was driven by the prompt structure and not simply thermal decoding noise \cite{wang2023self}. The Repeat Control CARs were exceptionally high: 0.913 for the 0.5B model and 0.813 for the 1.5B model. Because the prompt-variant agreement is substantially lower than the repeat-run agreement, we successfully isolate structural prompt instability as the primary failure vector \cite{mizrahi2024state}. 

Mean wording similarity across variants was extraordinarily low (0.202 and 0.166), indicating the models generated highly varied rationales \cite{kamalloo2023evaluating}. By evaluating CAR exclusively on the parsed logical claim, we decoupled superficial verbose differences from actual scientific contradictions \cite{elazar2021measuring}.

\begin{figure*}[htbp]
    \centering
    \includegraphics[width=0.8\textwidth]{1.2.png}
    \caption{Conclusion Agreement Rate (CAR) vs. Mean Wording Similarity across models. The stark disparity confirms that conclusion instability is distinct from mere lexical variance.}
    \label{fig:car_similarity}
\end{figure*}

\section{Discussion and Limitations}
The integration of LLMs into multi-step agentic research workflows assumes a foundational level of prompt stability \cite{wang2024survey, xi2023rise}. If an extraction agent has a 25\% to 50\% chance of flipping a conclusion based on minor formatting or positional variances, these errors will propagate \cite{wu2023agentbench}. The downstream verification agents, receiving hallucinated or flipped claims as their input state, will compound the error, leading to systemic workflow failure \cite{mialon2023augmented}. To mitigate this, developers of AI research tools must implement multi-prompt ensemble architectures rather than relying on a single, highly optimized prompt \cite{zhao2021calibrate}.

This empirical implementation is constrained by the sample size (30 questions) necessitated by single-GPU VRAM limits and local execution time constraints \cite{jin2019pubmedqa}. Calculating absolute failure rates with tight confidence intervals requires scaling to the full 1,000-question PubMedQA dataset. Additionally, the domain was strictly restricted to biomedical question-answering abstracts. 

\section{Conclusion}
This study provides empirical evidence that prompt variations cause severe semantic claim drift in Large Language Models \cite{errica2025wrong, sclar2024quantifying}. By utilizing a deterministic parser to extract structured claims, we demonstrated that formatting, positional, and semantic perturbations lead to high conclusion flip rates that cannot be attributed to sampling variance \cite{wang2023self}. 

While scaling model parameters from 0.5B to 1.5B significantly improves overall stability and mitigates catastrophic surface formatting failures, both models remain highly fragile \cite{qwen2024}. The 1.5B model's 50\% flip rate under semantic paraphrasing proves that a conclusion obtained from a single LLM prompt cannot be treated as a robust, fully tested scientific result \cite{turpin2023language}. 

\begin{thebibliography}{32}
\bibitem{errica2025wrong} F. Errica, G. Siracusano, D. Sanvito, and R. Bifulco, ``What did I do wrong? Quantifying LLMs' sensitivity and consistency to prompt engineering,'' \textit{Proceedings of the 2025 Conference of the Nations of the Americas Chapter of the Association for Computational Linguistics (NAACL 2025)}, vol. 1, pp. 1543--1558, 2025.
\bibitem{sclar2024quantifying} M. Sclar, Y. Choi, Y. Tsvetkov, and A. Suhr, ``Quantifying language models' sensitivity to spurious features in prompt design or: How I learned to start worrying about prompt formatting,'' \textit{Proceedings of the Twelfth International Conference on Learning Representations (ICLR 2024)}, 2024.
\bibitem{zheng2023judging} L. Zheng, W.-L. Chiang, Y. Sheng, S. Zhuang, Z. Wu, Y. Zhuang, Z. Lin, Z. Li, D. Li, E. P. Xing, H. Zhang, J. E. Gonzalez, and I. Stoica, ``Judging LLM-as-a-judge with MT-Bench and Chatbot Arena,'' \textit{Advances in Neural Information Processing Systems 36 (NeurIPS 2023), Datasets and Benchmarks Track}, 2023.
\bibitem{zhao2021calibrate} T. Z. Zhao, E. Wallace, S. Feng, D. Klein, and S. Singh, ``Calibrate before use: Improving few-shot performance of language models,'' \textit{Proceedings of the 38th International Conference on Machine Learning (ICML 2021)}, PMLR vol. 139, pp. 12697--12706, 2021.
\bibitem{lu2022fantastically} Y. Lu, M. Bartolo, A. Moore, S. Riedel, and P. Stenetorp, ``Fantastically ordered prompts and where to find them: Overcoming few-shot prompt order sensitivity,'' \textit{Proceedings of the 60th Annual Meeting of the Association for Computational Linguistics (ACL 2022)}, pp. 8086--8098, 2022.
\bibitem{mizrahi2024state} M. Mizrahi, G. Kaplan, D. Malkin, R. Dror, D. Shahaf, and G. Stanovsky, ``State of what art? A call for multi-prompt LLM evaluation,'' \textit{Transactions of the Association for Computational Linguistics}, vol. 12, pp. 933--949, 2024.
\bibitem{wang2023self} X. Wang, J. Wei, D. Schuurmans, Q. Le, E. Chi, S. Narang, A. Chowdhery, and D. Zhou, ``Self-consistency improves chain of thought reasoning in language models,'' \textit{Proceedings of the Eleventh International Conference on Learning Representations (ICLR 2023)}, 2023.
\bibitem{webson2022do} A. Webson and E. Pavlick, ``Do prompt-based models really understand the meaning of their prompts?,'' \textit{Proceedings of the 2022 Conference of the North American Chapter of the Association for Computational Linguistics (NAACL 2022)}, pp. 2300--2344, 2022.
\bibitem{turpin2023language} M. Turpin, J. Michael, E. Perez, and S. R. Bowman, ``Language models don't always say what they think: Unfaithful explanations in chain-of-thought prompting,'' \textit{Advances in Neural Information Processing Systems 36 (NeurIPS 2023)}, 2023.
\bibitem{reynolds2021prompt} L. Reynolds and K. McDonell, ``Prompt programming for large language models: Beyond the few-shot paradigm,'' \textit{arXiv preprint arXiv:2102.07350}, 2021.
\bibitem{kojima2022large} T. Kojima, S. S. Gu, M. Reid, Y. Matsuo, and Y. Iwasawa, ``Large language models are zero-shot reasoners,'' \textit{Advances in Neural Information Processing Systems 35}, pp. 22199--22213, 2022.
\bibitem{zhou2023large} Y. Zhou, A. I. Muresanu, Z. Han, K. Paster, S. Pitis, H. Chan, and J. Ba, ``Large language models are human-level prompt engineers,'' \textit{Proceedings of the Eleventh International Conference on Learning Representations (ICLR 2023)}, 2023.
\bibitem{jiang2020how} Z. Jiang, F. X. Wu, Z. Jiang, S. Wang, and G. Neubig, ``How can we know what language models know?,'' \textit{Transactions of the Association for Computational Linguistics}, vol. 8, pp. 423--438, 2020.
\bibitem{gao2021making} T. Gao, A. Fisch, and D. Chen, ``Making pre-trained language models better few-shot learners,'' \textit{Proceedings of the 59th Annual Meeting of the Association for Computational Linguistics (ACL 2021)}, pp. 3816--3830, 2021.
\bibitem{perez2021true} E. Perez, D. Kiela, and K. Cho, ``True few-shot learning with language models,'' \textit{Advances in Neural Information Processing Systems 34}, pp. 11054--11070, 2021.
\bibitem{schick2021exploiting} T. Schick and H. Schütze, ``Exploiting cloze-questions for few-shot text classification and natural language inference,'' \textit{Proceedings of the 16th Conference of the European Chapter of the Association for Computational Linguistics (EACL 2021)}, pp. 255--269, 2021.
\bibitem{brown2020language} T. B. Brown, B. Mann, N. Ryder, M. Subbiah, J. D. Kaplan, P. Dhariwal, A. Neelakantan, P. Shyam, G. Sastry, A. Askell et al., ``Language models are few-shot learners,'' \textit{Advances in Neural Information Processing Systems 33}, pp. 1877--1901, 2020.
\bibitem{ouyang2022training} L. Ouyang, J. Wu, X. Jiang, D. Almeida, C. L. Wainwright, P. Mishkin, C. Zhang, S. Agarwal, K. Slama, A. Ray et al., ``Training language models to follow instructions with human feedback,'' \textit{Advances in Neural Information Processing Systems 35}, pp. 27730--27744, 2022.
\bibitem{touvron2023llama} H. Touvron, T. Lavril, G. Izacard, X. Martinet, M. Lachaux, T. Lacroix, B. Rozière, N. Goyal, E. Hambro, F. Azhar et al., ``LLaMA: Open and efficient foundation language models,'' \textit{arXiv preprint arXiv:2302.13971}, 2023.
\bibitem{qwen2024} Qwen Team, ``Qwen2.5 technical report,'' \textit{arXiv preprint arXiv:2409.12918}, 2024.
\bibitem{jin2019pubmedqa} Q. Jin, B. Dhingra, Z. Liu, W. Cohen, and X. Lu, ``PubMedQA: A dataset for biomedical research question answering,'' \textit{Proceedings of the 2019 Conference on Empirical Methods in Natural Language Processing (EMNLP 2019)}, pp. 2566--2577, 2019.
\bibitem{honovich2022true} O. Honovich, S. Shaham, I. Bowman, and O. Levy, ``TRUE: Evaluating factual consistency in text generation,'' \textit{Proceedings of the 2022 Conference on Empirical Methods in Natural Language Processing (EMNLP 2022)}, pp. 3905--3920, 2022.
\bibitem{kamalloo2023evaluating} E. Kamalloo, N. Dziri, C. L. A. Clarke, and D. Rafiei, ``Evaluating language models for document retrieval,'' \textit{Proceedings of the 46th International ACM SIGIR Conference on Research and Development in Information Retrieval (SIGIR 2023)}, pp. 2341--2345, 2023.
\bibitem{ji2023survey} Z. Ji, N. Lee, R. Frieske, T. Yu, D. Su, Y. Xu, E. Ishii, Y. J. Bang, A. Madotto, and P. Fung, ``Survey of hallucination in natural language generation,'' \textit{ACM Computing Surveys}, vol. 55, no. 12, pp. 1--38, 2023.
\bibitem{huang2023towards} J. Huang and K. C. Chang, ``Towards reasoning in large language models: A survey,'' \textit{Proceedings of the 61st Annual Meeting of the Association for Computational Linguistics (ACL 2023)}, pp. 1045--1062, 2023.
\bibitem{mialon2023augmented} G. Mialon, R. Dessì, M. Lomeli, C. Nalmpantis, R. Pasunuru, R. Rao, S. Rozière, A. Zycode, H. Schwenk, and G. Synnaeve, ``Augmented language models: A survey,'' \textit{Transactions on Machine Learning Research}, 2023.
\bibitem{wu2023agentbench} X. Wu, J. Li, Y. Wang, J. Zhang, Y. Wang, and L. Wang, ``AgentBench: Evaluating LLMs as agents,'' \textit{Proceedings of the Twelfth International Conference on Learning Representations (ICLR 2024)}, 2024.
\bibitem{wang2024survey} L. Wang, C. Ma, X. Feng, Z. Zhang, H. Yang, J. Zhang, Z. Wu, W. Wang, J. Zhang, Z. Zou et al., ``A survey on large language model based autonomous agents,'' \textit{Frontiers of Computer Science}, vol. 18, no. 6, p. 186345, 2024.
\bibitem{xi2023rise} Z. Xi, W. Chen, X. Guo, W. He, Y. Ding, B. Hong, M. Zhang, J. Wang, S. Jin, E. Zhou et al., ``The rise and potential of large language model based agents: A survey,'' \textit{arXiv preprint arXiv:2303.01150}, 2023.
\bibitem{gao2023building} Y. Gao, Y. Xiong, X. Gao, K. Jia, J. Pan, Y. Bi, Y. Dai, J. Sun, and H. Wang, ``Building applications with LLMs: A comprehensive survey,'' \textit{arXiv preprint arXiv:2312.05862}, 2023.
\bibitem{salinas2024} A. Salinas and F. Morstatter, ``The unequal vulnerability of large language models to prompt perturbations,'' \textit{Proceedings of the Association for Computational Linguistics (ACL 2024)}, 2024.
\bibitem{elazar2021measuring} Y. Elazar, N. Kassner, S. Ravfogel, A. Ravichander, E. Hovy, Y. Goldberg, and I. Gurevych, ``Measuring and improving consistency in pretrained language models,'' \textit{Transactions of the Association for Computational Linguistics}, vol. 9, pp. 1012--1031, 2021.
\end{thebibliography}

\end{document}
